\section{Постановка и доступные вычисления}
\label{sec:model}
Пусть $Q\subset\mathbb R^d$ является непустым компактным выпуклым
множеством, $\operatorname{diam}Q\le D$, $D>0$. Функция $f$ выпукла
на $Q$ и дважды непрерывно дифференцируема в его окрестности. Предположим
\begin{equation}\label{eq:gradient_smoothness}
 \|\nabla^2f(x)\|\le L_0+L_1\|\nabla f(x)\|,\qquad x\in Q,
 \quad L_0,L_1\ge0.
\end{equation}
Для векторов используется евклидова норма, для матриц спектральная.
Минимизатор $x^*\in Q$ существует. Он нужен для записи гарантии
точности, но не подаётся на вход метода. Выражения $f(x)-f(x^*)$ и
$\|\nabla f(x)\|$ записываем полностью, без отдельных обозначений.

Доступны вычисление функции, вычисление градиента и линейная
минимизация по всему $Q$. Последняя возвращает
\[
 s\in\arg\min_{v\in Q}\langle\nabla f(x),v\rangle.
\]
Полный зазор Франка-Вульфа вычисляется по формуле
\begin{equation}\label{eq:gap}
 G(x)=\max_{v\in Q}\langle\nabla f(x),x-v\rangle
     =\langle\nabla f(x),x-s\rangle.
\end{equation}
Поскольку $x\in Q$, имеем $G(x)\ge0$. Из выпуклости следует
$f(x)-f(x^*)\le G(x)$, поэтому $f(x)-G(x)$ даёт нижнюю границу
неизвестного минимума \cite{jaggi}. Сам минимум при проверке остановки
вычислять не требуется.

В минимуме на множестве выполняется условие
\[
 \langle\nabla f(x^*),v-x^*\rangle\ge0\quad(v\in Q).
\]
Оно равносильно $G(x^*)=0$ и не требует $\nabla f(x^*)=0$.
Например, $f(x)=Ax+x^2/2$ на $[0,1]$, $A>0$, имеет минимум в нуле,
производную $A$ в этой точке и гессиан, равный единице.
Большой градиент и большая кривизна являются разными свойствами.

\begin{lemma}[Конечная граница кривизны по начальному градиенту]
\label{lem:global_curvature}
Из условия гладкости следует
\begin{equation}\label{eq:global_curvature}
 \|\nabla^2f(z)\|\le K_Q\quad(z\in Q),\qquad
 K_Q=(L_0+L_1\|\nabla f(x_0)\|)e^{L_1D}.
\end{equation}
В результатах ниже вместо $K_Q$ можно использовать любую более точную
равномерную верхнюю границу. Эта константа нужна для оценки числа
вычислений; алгоритму знать её не требуется.
\end{lemma}
\begin{proof}
Для единичного направления $e$ вдоль отрезка из $x_0$ внутри $Q$
основная теорема анализа даёт
\[
 \|\nabla f(x_0+te)\|\le\|\nabla f(x_0)\|
 +\int_0^t\{L_0+L_1\|\nabla f(x_0+se)\|\}\,ds.
\]
При $L_1>0$ неравенство Гронуолла ограничивает
$L_0+L_1\|\nabla f(x_0+te)\|$ величиной
$(L_0+L_1\|\nabla f(x_0)\|)e^{L_1t}$. Остаётся учесть $t\le D$.
При $L_1=0$ исходное условие сразу даёт границу $L_0$.
\end{proof}

Условие гладкости не содержит неизвестного минимума. Алгоритм ниже
также не требует константы регулярности относительно неизвестного
минимизатора или его удалённости от границы. Полный зазор ограничивает
глобальную ошибку благодаря выпуклости функции.
